% ==================== 二、模型假设、符号与前提 ====================
\section{模型假设与符号约定}

\subsection{模型假设}

本文的建模基于以下假设。为便于评阅，每条假设均给出理由。

\begin{enumerate}[label=\textbf{A\arabic*.},itemsep=4pt,topsep=4pt,leftmargin=*]
  \item \textbf{段内功率恒定。}
        附件提供的负载、光伏数据以 10 分钟为间隔，本文视其为该区间内的平均功率，
        段内功率恒定，故每段电量 $=$ 功率 $\times \Delta t$。
        \emph{理由}：题面只给出离散采样点，采用平均功率是对连续量的标准离散化处理，
        且与附件 1 的“某天不同时间的电价/负载”表述一致。

  \item \textbf{储能效率模型为单向恒定效率。}
        充电效率与放电效率均为 $0.9$（往返 $0.81$），不计及功率相关性、
        温度效应与循环老化衰减。
        \emph{理由}：题附录 1 明确给出“充放电效率为 $90\%$”，
        未提供效率曲线或衰减规律，故采用题给常值。
        题面未区分“单向”与“往返”，本文在 \S\ref{sec:eff-robust} 给出两口径的完整对照。

  \item \textbf{不计储能折旧与运维费用。}
        目标函数中只包含购电相关费用。
        \emph{理由}：题面要求最小化购电费用，未提及设备成本；
        引入未经给定的折旧率反而会带来主观偏差。

  \item \textbf{严格因果性（非前视）。}
        任何决策（包括预测器的构造、融合权重的估计、场景的提取与场景树的分支）
        均只使用决策时刻之前\textbf{已实现}的数据，不得读取当天及未来的实际值。
        \emph{理由}：这是滚动决策问题可行性的前提，也是所有对照实验公平比较的基础。
        本文在问题一至问题四\textbf{逐问强化}该原则：问题一为全信息已知；
        问题二要求预测只依赖历史；问题三进一步要求场景树的节点划分也只能依据
        届时新发布的预报（\S\ref{sec:q3-nonanticip}）；问题四再加电价的非前视。

  \item \textbf{电价按供电时段计价。}
        每段费用 $=$ 该段的电价 $\times$ 该段的购电量，
        而不是按“作出决策的时刻”统一计价。
        \emph{理由}：物理上电量的计量发生在用电时刻；
        这使问题三的四个发布时刻各自对应一组时段价格，
        避免了把所有未来时段折算到决策时刻价格的非常规设定。

  \item \textbf{储能充放电互斥（但不显式约束）。}
        物理上储能不能同时充放电。本文\textbf{不加入二元变量}，
        而是证明在本文的模型结构下最优解自动满足互斥（\S\ref{sec:q1-tight}），
        并在输出前逐段检查。
        \emph{理由}：见 \S\ref{sec:q1-tight} 的等价性证明。
\end{enumerate}

\noindent 题面已明确给定的计费与物理规则（计划购电照付、取消部分退还原价并另付
$50\%$ 违约费、增购按 $1.5$ 倍、紧急购电按 $5$ 倍、不外送电且允许弃光、
各问的终端电量要求）不属于假设，而是\textbf{建模的已知条件}，
其统一处理见 \S\ref{sec:conventions}；与各问模型直接相关的终端条件
在对应问题的模型建立中陈述。

\subsection{全局符号说明}
\label{sec:notation}

四问共用同一套物理量，故统一在模型假设后集中说明，各问不再重复定义。
表 \ref{tab:notation} 给出全文使用的符号。

\begin{table}[H]
\centering
\caption{全文符号说明（四问通用）}
\label{tab:notation}
\small
\begin{tabularx}{\textwidth}{@{}c l X c@{}}
\toprule
符号 & 含义 & 说明 & 单位 \\
\midrule
\multicolumn{4}{@{}l}{\textit{（一）时段与数据}}\\
\midrule
$T$            & 时段数                   & 一天划分为 $T=144$ 个十分钟时段 & — \\
$\Delta t$     & 每段时长                 & $\Delta t=1/6$ h & h \\
$t$            & 时段索引                 & $t=1,\dots,144$ & — \\
$d$            & 日期索引                 & 年度模型中的日 & — \\
$p_t$          & 第 $t$ 段电价             & 问题一至三用附件 1；问题四用附件 4 & 元/kWh \\
$L_t$          & 第 $t$ 段负载电量         & 附件中为功率，乘 $\Delta t$ 折算 & kWh \\
$V_t$          & 第 $t$ 段光伏发电量       & 同上 & kWh \\
$N_t$          & 净负载                   & $N_t=L_t-V_t$ & kWh \\
\midrule
\multicolumn{4}{@{}l}{\textit{（二）决策变量}}\\
\midrule
$g_t$          & 计划购电量               & 日前制定的购电量 & kWh \\
$g_t^{0}$      & 0 点原计划购电量          & 问题三、4-3 的初始计划 & kWh \\
$q_t$          & 最终调整后购电量          & 问题三、4-3 的执行量 & kWh \\
$c_t$          & 充电量                   & 作用于设备外端 & kWh \\
$d_t$          & 放电量                   & 作用于设备外端 & kWh \\
$w_t$          & 弃光量                   & 仅设非负约束，见 \S\ref{sec:q1-model} & kWh \\
$E_t$          & 段末储能电量             & 内部库存，$E_t\in[1200,\,10800]$ & kWh \\
$u_t,v_t$      & 相对原计划的增购量、减购量 & $u_t=(q_t-g_t^0)_+$，$v_t=(g_t^0-q_t)_+$ & kWh \\
$e_t$          & 紧急购电量               & 缺口补救，按 $5p_t$ 计价 & kWh \\
\midrule
\multicolumn{4}{@{}l}{\textit{（三）效率与储能参数}}\\
\midrule
$\eta_c,\eta_d$ & 充电效率、放电效率       & 主口径 $\eta_c=\eta_d=0.9$，往返 $0.81$ & — \\
$\rho$         & 往返效率                 & $\rho=\eta_c\eta_d$ & — \\
$E_0$          & 初始储电量               & 题给 2025-01-01 0:00 为 6000 & kWh \\
\midrule
\multicolumn{4}{@{}l}{\textit{（四）随机规划量}}\\
\midrule
$\hat N_t$     & 净负载点预测             & 只用历史数据构造 & kWh \\
$N_{s,t}$      & 场景净负载               & 场景 $s$ 的第 $t$ 段净负载 & kWh \\
$p_{s,t}$      & 场景电价                 & 仅问题四使用 & 元/kWh \\
$s$            & 场景索引                 & 由历史日整日残差构成 & — \\
$w_s$          & 场景权重                 & 等权时 $w_s=1/S$ & — \\
$n$            & 场景树节点索引           & 问题三、4-3 & — \\
$\pi_n$        & 节点概率                 & 落入该节点的历史日数占比 & — \\
$\alpha$       & 融合权重                 & 在线估计，见 \S\ref{sec:q3-fusion} & — \\
\midrule
\multicolumn{4}{@{}l}{\textit{（五）评价指标}}\\
\midrule
$J$            & 总费用                   & 分项构成见对应公式 & 元 \\
$\lambda_t$    & 边际电价（对偶变量）     & 能量平衡约束的影子价格 & 元/kWh \\
\bottomrule
\end{tabularx}
\end{table}

\noindent 记号说明：$(\cdot)_+=\max\{\cdot,0\}$；$\E[\cdot]$ 为数学期望；
带下标 $s$ 的量表示随场景变化、由对应的历史日产生。
为避免与日期索引 $d$ 混淆，全文用 $d_t$ 表示放电量（与主程序中的变量名一致）。
